\documentclass[11pt]{article}
\usepackage{preprint}

\title{Language as a Tool that Shapes Identity}
\author{Jiadi He\email{jh2288@georgetown.edu}}
\affil{Georgetown University}
\shorttitle{Language as a Tool that Shapes Identity}
\keywords{language; media; representation}

\begin{document}
\maketitle
\keywordsline

Language plays an important role in the way we see the world and the media is integral in the way we understand what is going on in the world around us. We are consciously or unconsciously seeking information from the media and allowing it to be what we consider facts and a lens through which we view the world. The way in which news is structured influences how it is understood and the concepts in seeks to reproduce. For instance, after 9/11, the media in the United States started an anti-Arab campaign that made the citizens of the United States into victims whereas the Arabs became villains. The result of this phrasing by the media was that there were discriminatory actions that were taken against all Arabs regardless of whether or not they were part of a terrorist group. The situation shows how easily the media can influence the way people think about certain important aspects. The two texts \textit{Introduction of Contagious- Cultures, Carriers, and the Outbreak Narrative} by Priscilla Wald and \textit{Assaults on the Ivory Tower- Representation of Madness in the Discourse} by Price Margaret both show the dangers of language used by the media and the general public in describing incidents and how they it can lead to discrimination and marginalization, aspects that have negative implications on members of a group.

I enjoyed reading both texts and I could not help but see the similarities they share in how they discuss the aspect of language and how it frames different narratives. The text by Wald establishes the way in which reporting can make villains out of ordinary people by placing blame on them for something that was beyond their control. For instance, she discusses how after the outbreak of SARS, people started placing blame on each other as they sought to understand the cause and implication of the disease. The “outbreak narrative” that was created worked to divide people by placing the blame for the disease on patients. The individuals through which the virus spread were being villainized and myths created to show their role in the deaths that were being caused by the pandemic (3). It is sad how they had been blamed and shamed for things that were beyond their control, but that is what happens when myths are used to define science. Wald explains how the narrative of a twenty-six-year-old flight attendant who had been infected was defined by stating she “turned into a national scapegoat when Singapore’s minister of health announced at a press conference in early April that she “infected the whole lot of us”” (3). The blame game worked to create unwarranted hatred toWalds her that made people identify her with the disease. I find the narrative sickening because during the outbreak of the Corona Virus, Asians specifically Chinese nationals were being looked at with fear because they were a threat when it came to spreading the virus as portrayed in the media. People were being discriminated against because a misplaced “outbreak narrative” that marginalized and villainized a particular individual and race.

Identity is created in the way we describe things or people and the language we use and how we use it works magic in emphasizing the narrative. Wald explains that the myths that are created about scientific facts maniform the public by creating a new reality that differs from facts. She explains that the myths as a “strong emotional appeal derives from and affirms the fundamental values, hierarchies, and taxonomies that are the preconditions of that identity” (9). Myths around scientific issues such as the spread of viruses only work to create fear and anxiety around the topic and misguided blame directed at entities that had little to do with the emergence or spread of the virus. Wald’s text explains how the use of myths have been adopted in some scientific journals and reporting resources and the unpleasant nature of its social implications. When reading through the text, I questioned the role of myth in science because it did not occur to me that myths were a part of the scientific world. Wald however, made it clear that over time, the interconnected worlds of social theorists and scientists has led to them borrowing from one another and that has influenced how individuals interpret issues like contagious diseases and the outbreak process.

Wald’s texts provided me with the lens to better understand what Price talks about in her text and the reporting around the two shooters. Price explains that the mental health of Cho and Kazmierczak, the shooters, was used to create a myth around individuals with mental health conditions, making them into villains in the eyes of the general population. The fact that their actions were blamed on or justified using the fact that they had mental health issues works to marginalize and label individuals with mental health conditions as being violent and dangerous. According to Price, “these representations reify the belief that madness and sanity are two extremely separate spaces—one dangerous, one safe” (145). The invisible line drawn between those that are considered sane and those that are mad, works to sanitize the sane from the violent acts that were committed by Cho and Kazmierczak, and shift all the blame to their mental capacity. I was able to better understand the dangers of language as put forWald by Wald when she explained how the flight attendant in her text was villainized because she was a carrier for a virus without her knowledge (Wald 3). Similar to the outbreak narrative, the mental health narrative is meant to put one group of people against another as a way to sanitize one group while blaming the other. It is so unfortunate that a whole community of people who were already marginalized, were put on the spotlight for such atrocious crimes. Mental health is something that is beginning to get the recognition it deserves and making it synonymous to violence and murder only works to destroy the safe space that protects individuals with mental health conditions.

Reading through the case studies that Price provides in her text, I realized just how biased and stereotypical the media can be and the role it plays in enabling the spread of narrative that are based on myths. The case studies work to construct a narrative that shows that mental disability of both individuals was the cause of the deviant behavior exhibited through the shootings (Price 49). Disability is used as a scapegoat on which the actions of Cho and Kazmierczak can be blamed without taking into account the impact it has on the members of the already marginalized community. The cases of both Cho and Kazmierczak are described using stereotypes that fuel the myths about the individuals and their background and works to persuade the readers of the articles to draw a connection. The outbreak narrative is similar to the cases of Cho and Kazmierczak since it shows what the cause of the problem was by probing at different aspects to try and find answers to the problem. The only issue is how in the process of finding the source, we end up creating myths and narratives that destroy instead of enable. In the process individuals and groups find themselves marginalized and that is saddening.

Reading the two texts has made me more conscious of the misleading nature of the media in the language it uses and the narratives it propagates. I have realized that media literacy is important and that everyone using media should have access to it. The framing of narratives to create a villain and a victim only works to derail us as human beings and I believe instead of blaming one another we should work together to find solutions to the hard realities of life like disease and mass shootings.

\begin{references}[Works Cited]
\item Price, Margaret. ``Assaults on the Ivory tower: Representations of Madness in the Discourse of US School Shootings.'' \textit{Mad at School: Rhetorics of Mental Disability and Academic Life} (2011): 141-75.
\item Wald, P. (2020). Contagious: The outbreak narrative.
\end{references}

\end{document}
